\ifdefined\gkver\else\def\gkver{student}\fi
\documentclass[\gkver]{gaokaozhenti}
\begin{document}

\chapter{2016年四川理综}
%% paperType: 理综物理部分

\begin{enumerate}
\item
%% number: 1
%% paperName: 2016年四川理综第 1 题 6 分
%% typeId: 1
%% type: 单选题
%% chapter: 功和能
%% point: 动能定理
%% method: 无
%% score: 6
%% degree: 750
%% duplicateId: 0
%% body:
韩晓鹏是我国首位在冬奥会雪上项目夺冠的运动员。他在一次自由式滑雪空中技巧比赛中沿“助滑区”保持同一姿态下滑了一段距离，重力对他做功 $1900\UJ$，他克服阻力做功 $100\UJ$。韩晓鹏在此过程中\xzanswer{C}
\fourchoices[answer=C]
{动能增加了 $1900\UJ$}
{动能增加了 $2000\UJ$}
{重力势能减小了 $1900\UJ$}
{机械能减小了 $2000\UJ$}
%% answer: C
%% memo:
\memoanswer{
本题原卷及材料中未提供详解，待补充。
}

\item
%% number: 2
%% paperName: 2016年四川理综第 2 题 6 分
%% typeId: 1
%% type: 单选题
%% chapter: 交流电
%% point: 变压器
%% method: 无
%% score: 6
%% degree: 750
%% duplicateId: 0
%% body:
如图所示，接在家庭电路上的理想降压变压器给小灯泡 L 供电，如果将原、副线圈减少相同匝数，其它条件不变，则\xzanswer{B}
\onepicture[w=5.0cm]{figs/02.png}
\fourchoices[answer=B]
{小灯泡变亮}
{小灯泡变暗}
{原、副线圈两段电压的比值不变}
{通过原、副线圈电流的比值不变}
%% answer: B
%% memo:
\memoanswer{
根据变压器电压与匝数关系，$\frac{U_{1}}{U_{2}}=\frac{n_{1}}{n_{2}}$，因为是降压变压器，则 $n_{1}>n_{2}$，则当原、副线圈减少相同匝数时，由数学知识可知 $\frac{n_{1}}{n_{2}}$ 变大，则 $U_{2}$ 减小，故灯泡变暗，选项AC错误，B正确；根据 $\frac{I_{1}}{I_{2}}=\frac{n_{2}}{n_{1}}$ 可知通过原、副线圈电流的比值变小，选项D错误。
}

\item
%% number: 3
%% paperName: 2016年四川理综第 3 题 6 分
%% typeId: 1
%% type: 单选题
%% chapter: 曲线运动
%% point: 人造地球卫星
%% method: 无
%% score: 6
%% degree: 450
%% duplicateId: 0
%% body:
国务院批复，自2016年起将4月24日设立为“中国航天日”。1970年4月24日我国首次成功发射的人造卫星东方红一号，目前仍然在椭圆轨道上运行，其轨道近地点高度约为 $440\Ukm$，远地点高度约为 $2060\Ukm$；1984年4月8日成功发射的东方红二号卫星运行在赤道上空 $35786\Ukm$ 的地球同步轨道上。设东方红一号在远地点的加速度为 $a_{1}$，东方红二号的加速度为 $a_{2}$，固定在地球赤道上的物体随地球自转的加速度为 $a_{3}$，则 $a_{1}$、$a_{2}$、$a_{3}$ 的大小关系为\xzanswer{D}
\onepicture[w=5.0cm]{figs/03.png}
\fourchoices[answer=D]
{$a_{2}>a_{1}>a_{3}$}
{$a_{3}>a_{2}>a_{1}$}
{$a_{3}>a_{1}>a_{2}$}
{$a_{1}>a_{2}>a_{3}$}
%% answer: D
%% memo:
\memoanswer{
本题原卷及材料中未提供详解，待补充。
}

\item
%% number: 4
%% paperName: 2016年四川理综第 4 题 6 分
%% typeId: 1
%% type: 单选题
%% chapter: 磁场
%% point: 带电粒子在磁场中的运动
%% method: 无
%% score: 6
%% degree: 550
%% duplicateId: 0
%% body:
如图所示，正六边形abcdef区域内有垂直于纸面的匀强磁场。一带正电的粒子从 $f$ 点沿fd方向射入磁场区域，当速度大小为 $v_{b}$ 时，从b点离开磁场，在磁场中运动的时间为 $t_{b}$，当速度大小为 $v_{c}$ 时，从c点离开磁场，在磁场中运动的时间为 $t_{c}$，不计粒子重力。则\xzanswer{A}
\onepicture[w=4.0cm]{figs/04.png}
\fourchoices[answer=A]
{$v_{b}:v_{c}=1:2$，$t_{b}:t_{c}=2:1$}
{$v_{b}:v_{c}=2:1$，$t_{b}:t_{c}=1:2$}
{$v_{b}:v_{c}=2:1$，$t_{b}:t_{c}=2:1$}
{$v_{b}:v_{c}=1:2$，$t_{b}:t_{c}=1:2$}
%% answer: A
%% memo:
\memoanswer{
本题原卷及材料中未提供详解，待补充。
}

\item
%% number: 5
%% paperName: 2016年四川理综第 5 题 6 分
%% typeId: 1
%% type: 单选题
%% chapter: 光学
%% point: 光的折射
%% method: 无
%% score: 6
%% degree: 450
%% duplicateId: 0
%% body:
某同学通过实验测定半圆形玻璃砖的折射率 $n$。如图\ref{2016四川理综05a}所示，O是圆心，MN是法线，AO、BO分别表示某次测量时光线在空气和玻璃砖中的传播路径。该同学测得多组入射角 $i$ 和折射角 $r$，做出 $\sin i-\sin r$ 图像如图\ref{2016四川理综05b}所示。则\xzanswer{B}
\twopicture[labelA=2016四川理综05a, labelB=2016四川理综05b, numA=甲, numB=乙, wA=5.0cm, wB=5.0cm]
{figs/05a.png}
{figs/05b.png}
\fourchoices[answer=B]
{光由A经O到B，$n=1.5$}
{光由B经O到A，$n=1.5$}
{光由A经O到B，$n=0.67$}
{光由B经O到A，$n=0.67$}
%% answer: B
%% memo:
\memoanswer{
由图线可知 $\frac{\sin i}{\sin r}=\frac{0.6}{0.9}=\frac{2}{3}=\frac{1}{n}$，可得 $n=1.5$；因 $i$ 是入射角，$r$ 是折射角，折射角大于入射角，故光由B经O到A，故选B。
}

\item
%% number: 6
%% paperName: 2016年四川理综第 6 题 6 分
%% typeId: 2
%% type: 多选题
%% chapter: 机械振动
%% point: 振动图象
%% method: 无
%% score: 6
%% degree: 450
%% duplicateId: 0
%% body:
简谐横波在均匀介质中沿直线传播，P、Q是传播方向上相距 $10\Um$ 的两质点，波先传到P，当波传到Q开始计时，P、Q两质点的振动图像如图所示。则\xzanswer{AD}
\onepicture[w=5.5cm]{figs/06.png}
\fourchoices[answer=AD]
{质点 Q 开始振动的方向沿 y 轴正方向}
{该波从P传到Q的时间可能为 $7\Us$}
{该波的传播速度可能为 $2\Ums$}
{该波的波长可能为 $6\Um$}
%% answer: AD
%% memo:
\memoanswer{
本题原卷及材料中未提供详解，待补充。
}

\item
%% number: 7
%% paperName: 2016年四川理综第 7 题 6 分
%% typeId: 2
%% type: 多选题
%% chapter: 电磁感应
%% point: 电磁感应中图象
%% method: 无
%% score: 6
%% degree: 350
%% duplicateId: 0
%% body:
如图所示，电阻不计、间距为 $l$ 的光滑平行金属导轨水平放置于磁感应强度为 $B$、方向竖直向下的匀强磁场中，导轨左端接一定值电阻 $R$。质量为 $m$、电阻为 $r$ 的金属棒MN置于导轨上，受到垂直于金属棒的水平外力 $F$ 的作用由静止开始运动，外力 $F$ 与金属棒速度 $v$ 的关系是 $F=F_{0}+kv$（$F_{0}$、$k$ 是常量），金属棒与导轨始终垂直且接触良好。金属棒中感应电流为 $i$，受到的安培力大小为 $F_{A}$，电阻 $R$ 两端的电压为 $U_{R}$，感应电流的功率为 $P$，它们随时间 $t$ 变化图像可能正确的有\xzanswer{BC}
\onepicture[w=4.5cm]{figs/07.png}
\fourchoices[answer=BC, ispicture=true, h=2.4cm]
{figs/07a.png}
{figs/07b.png}
{figs/07c.png}
{figs/07d.png}
%% answer: BC
%% memo:
\memoanswer{
根据牛顿定理可知，某时刻金属棒的加速度为 $a=\frac{F-F_{A}}{m}=\frac{F_{0}+kv-\frac{B^{2}L^{2}v}{R+r}}{m}$。

若 $k>\frac{B^{2}L^{2}}{R+r}$，则金属棒做加速度增加的加速运动，其 $v-t$ 图象如图1所示；导体的电流 $I=\frac{BLv}{R+r}$ 可知 $I$ 与 $v$ 成正比，则 $I-t$ 图线应该与 $v-t$ 图线形状相同。根据 $F_{A}=\frac{B^{2}L^{2}v}{R+r}$ 可知，$F_{A}$ 与 $v$ 成正比，则 $F_{A}-t$ 图线应该和 $v-t$ 图线形状相同，选项B正确。

\twopicture[showlabel=false, wA=5.5cm, wB=5.5cm]
{figs/07_an1.png}
{figs/07_an2.png}

根据 $P=\frac{E^{2}}{R+r}=\frac{B^{2}L^{2}v^{2}}{R+r}$ 可知 $P$ 与 $v^{2}$ 成正比，则 $P-t$ 图线不应该是直线，选项D错误。

同理，若 $k<\frac{B^{2}L^{2}}{R+r}$，金属棒做加速度减少的加速运动，其 $v-t$ 图像如图2所示，$I-t$ 图线应该和 $v-t$ 图线形状相同；若 $k=\frac{B^{2}L^{2}}{R+r}$，金属棒做匀加速直线运动，$I-t$ 因为过原点的倾斜直线，不会保持不变。选项A错误。

根据 $U_{R}=\frac{BLvR}{R+r}$ 可知 $U_{R}$ 与 $v$ 成正比，则 $U_{R}-t$ 图线应该和 $v-t$ 图线形状相同，选项C正确。

此题的正确选项为 BC。
}

\item
%% number: 8
%% paperName: 2016年四川理综第 8 题 11 分
%% typeId: 4
%% type: 实验题
%% chapter: 电路
%% point: 测电动势与内阻
%% method: 无
%% score: 11
%% degree: 450
%% duplicateId: 0
%% body:
\textbf{Ⅰ．}（6 分）用如图所示的装置测量弹簧的弹性势能。将弹簧放置在水平气垫导轨上，左端固定，右端在O点；在O点右侧的B、C位置各安装一个光电门，计时器（图中未画出）与两个光电门相连。先用米尺测得B、C两点间距离 $s$，再用带有遮光片的滑块压缩弹簧到某位置A，静止释放，计时器显示遮光片从B到C所用的时间 $t$，用米尺测量A、O之间的距离 $x$。
\begin{enumerate}
	\item
	计算滑块离开弹簧时速度大小的表达式是 \tkanswer[2cm]{$v=\frac{s}{t}$}。
	\item
	为求出弹簧的弹性势能，还需要测量 \tkanswer[1.5cm]{C}。

	A．弹簧原长\quad B．当地重力加速度\quad C．滑块（含遮光片）的质量
	\item
	增大A、O之间的距离 $x$，计时器显示时间 $t$ 将 \tkanswer[1.5cm]{B}。

	A．增大\quad B．减小\quad C．不变
\end{enumerate}
\onepicture[w=7.5cm, align=right]{figs/08a.png}

\textbf{Ⅱ．}（11 分）用如图所示电路测量电源的电动势和内阻。实验器材：待测电源（电动势约 $3\UV$，内阻约 $2\UO$），保护电阻 $R_{1}$（阻值 $10\UO$）和 $R_{2}$（阻值 $5\UO$），滑动变阻器 $R$，电流表A，电压表V，开关S，导线若干。

实验主要步骤：
\begin{steps}
	\item
	将滑动变阻器接入电路的阻值调到最大，闭合开关；
	\item
	逐渐减小滑动变阻器接入电路的阻值，记下电压表的示数 $U$ 和相应电流表的示数 $I$；
	\item
	以 $U$ 为纵坐标，$I$ 为横坐标，作 $U-I$ 图线（$U$、$I$ 都用国际单位）；
	\item
	求出 $U-I$ 图线斜率的绝对值 $k$ 和在横轴上的截距 $a$。
\end{steps}
回答下列问题：
\begin{enumerate}
	\item
	电压表最好选用 \tkanswer[1.5cm]{A}；电流表最好选用 \tkanswer[1.5cm]{C}。

	A．电压表（$0\sim3\UV$，内阻约 $15\UkO$）\quad B．电压表（$0\sim3\UV$，内阻约 $3\UkO$）
	C．电流表（$0\sim200\UmA$，内阻约 $2\UO$）\quad D．电流表（$0\sim30\UmA$，内阻约 $2\UO$）
	\item
	滑动变阻器的滑片从左向右滑动，发现电压表示数增大。两导线与滑动变阻器接线柱连接情况是 \tkanswer[1.5cm]{C}。

	A．两导线接在滑动变阻器电阻丝两端接线柱
	B．两导线接在滑动变阻器金属杆两端接线柱
	C．一条导线接在滑动变阻器金属杆左端接线柱，另一条导线接在电阻丝左端的接线柱
	D．一条导线接在滑动变阻器金属杆右端接线柱，另一条导线接在电阻丝右端的接线柱
	\item
	选用 $k$、$a$、$R_{1}$ 和 $R_{2}$ 表示待测电源的电动势 $E$ 和内阻 $r$ 的表达式 $E=$ \tkanswer[2cm]{$ka$}，$r=$ \tkanswer[2cm]{$k-R_{2}$}，代入数值可得 $E$ 和 $r$ 的测量值。
\end{enumerate}
\onepicture[w=6.0cm, align=right]{figs/08b.png}
%% answer: 
\jdanswer{
\begin{enumerate}
	\item
	Ⅰ．（1）$v=\frac{s}{t}$；（2）C；（3）B。

	\item
	Ⅱ．（1）A、C；（2）C；（3）$ka$；$k-R_{2}$。

\end{enumerate}
}
%% memo:
\memoanswer{
本题原卷及材料中未提供详解，待补充。
}

\item
%% number: 9
%% paperName: 2016年四川理综第 9 题 15 分
%% typeId: 6
%% type: 计算题
%% chapter: 电场
%% point: 带电粒子的直线运动
%% method: 无
%% score: 15
%% degree: 650
%% duplicateId: 0
%% body:
中国科学院2015年10月宣布中国将在2020年开始建造世界上最大的粒子加速器。加速器是人类揭示物质本源的关键设备，在放射治疗、食品安全、材料科学等方面有广泛应用。

如图所示，某直线加速器由沿轴线分布的一系列金属圆管（漂移管）组成，相邻漂移管分别接在高频脉冲电源的两极。质子从K点沿轴线进入加速器并依次向右穿过各漂移管，在漂移管内做匀速直线运动，在漂移管间被电场加速，加速电压视为不变。设质子进入漂移管B时速度为 $8\times10^{6}\Ums$，进入漂移管E时速度为 $1\times10^{7}\Ums$，电源频率为 $1\times10^{7}\UHz$，漂移管间缝隙很小，质子在每个管内运动时间视为电源周期的 $\frac{1}{2}$。质子的荷质比取 $1\times10^{8}\UCkg$。求：
\begin{enumerate}
	\item
	漂移管 B 的长度；
	\item
	相邻漂移管间的加速电压。
\end{enumerate}
\onepicture[w=7.0cm, align=right]{figs/09.png}
%% answer: 
\jdanswer{
\begin{enumerate}
	\item
	$0.4\Um$

	\item
	$6\times10^{4}\UV$

\end{enumerate}
}
%% memo:
\memoanswer{
本题原卷及材料中未提供详解，待补充。
}

\item
%% number: 10
%% paperName: 2016年四川理综第 10 题 17 分
%% typeId: 6
%% type: 计算题
%% chapter: 运动和力
%% point: 牛顿定律的应用
%% method: 无
%% score: 17
%% degree: 450
%% duplicateId: 0
%% body:
避险车道是避免恶性交通事故的重要设施，由制动坡床和防撞设施等组成，如图竖直平面内，制动坡床视为与水平面夹角为 $\theta$ 的斜面。一辆长 $12\Um$ 的载有货物的货车因刹车失灵从干道驶入制动坡床，当车速为 $23\Ums$ 时，车尾位于制动坡床的底端，货物开始在车厢内向车头滑动，当货物在车厢内滑动了 $4\Um$ 时，车头距制动坡床顶端 $38\Um$，再过一段时间，货车停止。已知货车质量是货物质量的4倍，货物与车厢间的动摩擦因数为0.4；货车在制动坡床上运动受到的坡床阻力大小为货车和货物总重的0.44倍。货物与货车分别视为小滑块和平板，取 $\cos\theta=1$，$\sin\theta=0.1$，$g=10\Umsq$。求：
\begin{enumerate}
	\item
	货物在车厢内滑动时加速度的大小和方向；
	\item
	制动坡床的长度。
\end{enumerate}
\onepicture[w=9.0cm, align=right]{figs/10.png}
%% answer: 
\jdanswer{
\begin{enumerate}
	\item
	$5\Umsq$，方向沿斜面向下

	\item
	$98\Um$

\end{enumerate}
}
%% memo:
\memoanswer{
本题原卷及材料中未提供详解，待补充。
}

\item
%% number: 11
%% paperName: 2016年四川理综第 11 题 19 分
%% typeId: 6
%% type: 计算题
%% chapter: 磁场
%% point: 带电粒子在复合场中的运动
%% method: 无
%% score: 19
%% degree: 350
%% duplicateId: 0
%% body:
如图所示，图面内有竖直线 $DD'$，过 $DD'$ 且垂直于图面的平面将空间分成Ⅰ、Ⅱ两区域。区域Ⅰ有方向竖直向上的匀强电场和方向垂直图面的匀强磁场 $B$（图中未画出）；区域Ⅱ有固定在水平面上高 $h=2l$、倾角 $\alpha=\frac{\pi}{4}$ 的光滑绝缘斜面，斜面顶端与直线 $DD'$ 距离 $s=4l$，区域Ⅱ可加竖直方向的大小不同的匀强电场（图中未画出）；C点在 $DD'$ 上，距地面高 $H=3l$。零时刻，质量为 $m$、带电量为 $q$ 的小球P在K点具有大小 $v_{0}=\sqrt{gl}$、方向与水平面夹角 $\theta=\frac{\pi}{3}$ 的速度。在区域Ⅰ内做半径 $r=\frac{3l}{\pi}$ 的匀速圆周运动，经C点水平进入区域Ⅱ。某时刻，不带电的绝缘小球A由斜面顶端静止释放，在某处与刚运动到斜面的小球P相遇。小球视为质点，不计空气阻力及小球P所带电量对空间电磁场的影响。$l$ 已知，$g$ 为重力加速度。
\begin{enumerate}
	\item
	求匀强磁场的磁感应强度 $B$ 的大小；
	\item
	若小球A、P在斜面底端相遇，求释放小球A的时刻 $t_{A}$；
	\item
	若小球A、P在时刻 $t=\beta\sqrt{\frac{l}{g}}$（$\beta$ 为常数）相遇于斜面某处，求此情况下区域Ⅱ的匀强电场的场强 $E$，并讨论场强 $E$ 的极大值和极小值及相应的方向。
\end{enumerate}
\onepicture[w=7.0cm, align=right]{figs/11.png}
%% answer: 
\jdanswer{
\begin{enumerate}
	\item
	$B=\frac{\pi m}{3q}\sqrt{\frac{g}{l}}$；

	\item
	$t_{A}=(3-2\sqrt{2})\sqrt{\frac{l}{g}}$

	\item
	场强极小值为 $E_{min}=0$；场强极大值为 $E_{max}=\frac{7mg}{8q}$，方向竖直向上。

\end{enumerate}
}
%% memo:
\memoanswer{
本题原卷及材料中未提供详解，待补充。
}

\end{enumerate}

\end{document}
